\documentclass[10pt,a4paper,fontset=fandol]{ctexart}

\usepackage[a4paper,margin=1.3cm]{geometry}
\usepackage{amsmath}
\usepackage{booktabs}
\usepackage{graphicx}
\usepackage{caption}
\usepackage{enumitem}
\usepackage{mwe}
\usepackage[colorlinks=true,urlcolor=blue]{hyperref}

\setlength{\parindent}{0pt}
\setlength{\parskip}{2pt}
\setlist{nosep,leftmargin=1.5em}
\setlength{\abovedisplayskip}{3pt}
\setlength{\belowdisplayskip}{3pt}
\captionsetup{font=small,skip=2pt}

\ctexset{
  section={
    format=\large\bfseries,
    beforeskip=3pt,
    afterskip=2pt
  },
  subsection={
    format=\normalsize\bfseries,
    beforeskip=2pt,
    afterskip=1pt
  }
}

\title{\vspace{-2em}LaTeX 基本功能展示}
\author{徐宜安\quad 25060012033}
\date{2026 年 8 月 31 日}

\begin{document}

\maketitle
\vspace{-2.5em}

\begin{minipage}[t]{0.48\textwidth}

\section{文档结构}

\subsection{列表}

无序列表：
\begin{itemize}
  \item Shell 命令
  \item Git 版本控制
  \item LaTeX 排版
\end{itemize}

有序列表：
\begin{enumerate}
  \item 编写源文件
  \item 编译并检查 PDF
\end{enumerate}

\subsection{数学公式}

行内公式示例：$E=mc^2$。

带编号公式：
\begin{equation}
  S = 0.4A + 0.6B
  \label{eq:score}
\end{equation}

公式~\eqref{eq:score}表示加权成绩。

多行对齐：
\begin{align}
  x+y &= 10,\\
  x-y &= 2.
\end{align}

矩阵：
\[
A=
\begin{pmatrix}
1 & 2\\
3 & 4
\end{pmatrix}
\]

分段函数：
\[
f(x)=
\begin{cases}
x, & x\ge 0,\\
-x, & x<0.
\end{cases}
\]

\end{minipage}
\hfill
\begin{minipage}[t]{0.48\textwidth}

\section{图表与引用}

表~\ref{tab:plan}展示学习安排。

\begin{center}
\captionof{table}{学习安排}
\label{tab:plan}
\begin{tabular}{ccc}
\toprule
工具 & 时间 & 状态\\
\midrule
Shell & 2 h & 完成\\
Git & 3 h & 完成\\
LaTeX & 2 h & 进行中\\
\bottomrule
\end{tabular}
\end{center}

图~\ref{fig:demo}展示图片插入效果。

\begin{center}
  \includegraphics[width=0.55\linewidth]{example-image-a}
  \captionof{figure}{图片示例}
  \label{fig:demo}
\end{center}

\section{链接与文献}

课程参考资料：

\href{https://missing-semester-cn.github.io/}
{Missing Semester 中文版}

正文引用参考资料~\cite{missingsemester}。

\bibliographystyle{unsrt}
{\small\bibliography{references}}

\end{minipage}

\end{document}